\documentclass[../main.tex]{subfiles}
\begin{document}

\chapter{メモリオーダリング}
\lessoninfo{
  video={Lesson 9，動画 1--14},
  textbook={H\&P \cite{hennessy2017} §3.5--3.6},
  keywords={メモリを介した依存，ロード，ストア，コミット，ロードストアキュー，LSQ，ストアロードフォワーディング，ロードストアのアウトオブオーダ実行，ロードストアのインオーダ実行，再実行，リオーダバッファ，リザベーションステーション},
}

第8章のリオーダバッファによって，アウトオブオーダプロセッサはレジスタを
介するあらゆる依存を処理し，例外を正確に保てるようになった．しかしロード
とストアは，メモリを介しても互いに依存する．あるロードが，先行するストア
の書き込む位置を読むかどうかが分かるのは，両方のアドレスが計算された後で
ある．第7章がロードとストアをプログラム順に実行していたのはこのためで
ある．この章では，それらをアウトオブオーダに実行する方法を示す．ストアは
コミット時にはじめてメモリに書き込み，ロードストアキューが未コミットの
ロードとストアをプログラム順に保持し，各ロードは自分の値をメモリから，
あるいはそれを書いたストアから直接得る．

\section{メモリを介した依存}
\label{sec:l09-dependences}

第8章のプロセッサは，制御依存と，レジスタを介するあらゆる依存に対処する．
\source{Lesson 9，動画 1--2；H\&P §3.5--3.6．}分岐予測（第4章）は制御依存
のコストの大半を取り除き，ROB は予測を外した経路の命令を取り消す．
リネーミング（第6章）は，同じレジスタへの書き込みの間の偽の依存を除去
する．Tomasulo のアルゴリズムによるスケジューリング（第7章）は，各命令
を，それが読むレジスタ値が計算されるまで待たせる．そして ROB は正確な
例外を実現する．

残るのはロードとストアである．\margin{典型的な x86 プログラムが実行する
命令のおよそ 38\% はロードかストアであり，ロードはストアのおよそ2倍で
ある（Intel，\emph{Smart Memory Access} ホワイトペーパー，2006）．}これらは，
命令がレジスタを介して互いに依存するのとまったく同じように，メモリ位置を
介して互いに依存する
（\cref{sec:l07-load-store}）．
\begin{lstlisting}[language={[x86masm]Assembler}, numbers=none]
SW  R1, 0(R2)      ; MEM[R2] <- R1
ADD R4, R5, R6     ; R4 <- R5 + R6
LW  R3, 0(R4)      ; R3 <- MEM[R4]
\end{lstlisting}
R4 が R2 に等しいとき，かつそのときに限り，ロードはストアに対して RAW
依存を持つ．\caution{メモリにも WAR 依存と WAW 依存はあるが，ストアは
プログラム順にメモリへ書き込む（\cref{sec:l09-write}）ので，仕掛けが必要
なのは RAW だけである．}このような依存をレジスタを介した依存と区別する
性質が2つある．
\begin{itemize}
  \item 命令そのものからは見えない．レジスタ名はメモリ位置について何も
    語らず，メモリ位置はアドレスが計算されてはじめて，多くの場合その命令
    が発行された数サイクル後に分かる．
  \item メモリをリネーミングするものは存在しない．RAT が対応付けるのは
    レジスタ名だけなので，プロセッサはロードごとに，その値を読むべき
    ストアを見つけるか，そのようなストアがないと判定しなければならない．
\end{itemize}
この章の残りでは2つの問いに答える．ストアが実際にメモリに書き込むのは
いつか，そしてロードはどこで自分の値を見つけるか，である．

\section{メモリが書かれるとき}
\label{sec:l09-write}

ストアも他の命令と同様にオペランドが揃えば実行されるが，その時点では
メモリに書き込まない．\source{Lesson 9，動画 3；H\&P §3.6．}ストアの実行
はアドレスを計算し書き込む値を確定するだけであり，書き込み自体はストアが
コミットするときに行われる．メモリへの書き込みは，レジスタへの書き込みと
同様に取り消せない．ストアが実行時にメモリに書き込むとすると，その後に
先行命令が例外を発生させたとき，ハンドラは，プログラム順ではまだ書かれて
いないはずの位置を見ることになり，また予測を外した経路のストアは，
プログラムが変更しないメモリを変更してしまう．これに対しコミット時には，
そのストアがプログラムの実行に属することが確定しており，しかもストアは
プログラム順にコミットする．\sidenote{コミットとは書き込みが取り消せなく
なる時点であって，キャッシュに届く時点ではない．実際のコアでは，値はまず
ストアバッファに移り，その後 L1 キャッシュへ流し出される．他のコアがその
値を遅れて見ることはあり，それを制約するのが第21章の一貫性モデルである．}

書き込みを遅らせることは，ロードにとって問題を生む．ロードはアドレスを
計算し次第メモリを読むことができるが，プログラム順で先行するストアが，
実行を終えてコミットを待っている，あるいはまだアドレスすら計算していない
ことがある．そのストアがロードの読む位置に書き込むなら，メモリはまだ古い
値を保持しており，ロードは代わりにそのストアの値を得なければならない．
したがってプロセッサは，まだコミットしていないすべてのストア
のアドレスと値をプログラム順に保持し，すべてのロードがそれらを参照し
なければならない．\margin{第3章のパイプラインでは，MEM ステージがロードと
ストアをプログラム順に行う．}

\begin{keypoint}
  ストアはアウトオブオーダに実行されるが，メモリへの書き込みはコミット時
  にプログラム順で行う．それまでその値はプロセッサの内部にしか存在せず，
  ロードはそこを探さなければならない．
\end{keypoint}

\section{ロードストアキュー}
\label{sec:l09-lsq}

まだコミットしていないロードとストアは，ROB と同じように働くキューに
保持される．\source{Lesson 9，動画 4--5；H\&P §3.6．}

\begin{definition}[ロードストアキュー]
  \label{def:l09-lsq}
  \term[ろーどすとあきゅー]{ロードストアキュー}[load-store queue]%
  \index{LSQ|see{ロードストアキュー}}（LSQ）は，発行済みでまだコミットして
  いないすべてのロードとストアを，プログラム順に保持する．各エントリは
  4つのフィールドを持つ．
  \begin{description}
    \item[L/S] その命令がロードかストアか．
    \item[Addr] メモリアドレス．計算された後に記録される．
    \item[Value] ストアでは書き込む値，ロードでは得た値．
    \item[C] エントリが完了したときに立つ．ロードでは値を得たとき，
      ストアではアドレスと値の両方を記録したときである．
  \end{description}
\end{definition}

ROB と同様に，LSQ は発行ポインタとコミットポインタを持つ循環バッファで
ある（\cref{fig:l09-lsq}）．ロードやストアは，発行されるときに ROB
エントリに加えて発行ポインタの位置のエントリを受け取り，その LSQ エントリ
はコミット時に解放される．他の命令は LSQ を使わない．ロードはアドレスを
記録し，続いて値を記録する．ストアは，ベースレジスタが得られ次第アドレス
を，値レジスタが得られ次第値を，どちらの順でも記録する．ストアはコミット
するまでエントリを保持し，そのときにはじめてメモリに書き込むので，LSQ に
あるストアは，値がまだメモリに入っていないストアにちょうど一致する．
\margin{ROB もロードとストアを保持するが，キューを分けることでアドレスの
探索を短く保てる．}

\begin{figure}[htbp]
  \centering
  % The load-store queue: loads and stores in program order between the commit
% pointer and the issue pointer.  The load in LSQ6 searches the earlier
% entries, skips loads and non-matching stores, and takes the value of the
% nearest store with the same address.
\begin{tikzpicture}[x=1cm,y=1cm, font=\footnotesize\sffamily, >={Stealth[length=1.8mm]},
  hd/.style={font=\footnotesize\bfseries\sffamily},
  used/.style={draw, fill=ln-accent!12, minimum height=0.45cm, inner sep=0pt},
  free/.style={draw, fill=white, minimum height=0.45cm, inner sep=0pt},
  me/.style={draw, fill=ln-accent!35, minimum height=0.45cm, inner sep=0pt},
  lab/.style={font=\scriptsize\sffamily, text=ln-muted},
  ann/.style={font=\scriptsize\sffamily, anchor=west}]
  % column headings
  \node[hd] at (0.35,0.45) {L/S};
  \node[hd] at (1.25,0.45) {Addr};
  \node[hd] at (2.25,0.45) {Value};
  \node[hd] at (2.95,0.45) {C};
  % rows: entry / type / address / value / style
  \foreach \r/\t/\a/\v/\st in {1/{}/{}/{}/free, 2/S/0x100/17/used,
      3/L/0x108/4/used, 4/S/0x108/25/used, 5/L/{}/{}/used,
      6/L/0x100/{}/me, 7/S/{}/{}/used, 8/{}/{}/{}/free} {
    \pgfmathsetmacro{\y}{-(\r-1)*0.45}
    \node[\st, minimum width=0.7cm] at (0.35,\y) {\t};
    \node[\st, minimum width=1.1cm, font=\scriptsize\ttfamily] at (1.25,\y) {\a};
    \node[\st, minimum width=0.9cm] at (2.25,\y) {\v};
    \node[\st, minimum width=0.5cm] at (2.95,\y) {};
    \node[lab, anchor=east] at (-0.05,\y) {LSQ\r};
  }
  % completion bits: set for LSQ2-LSQ4, and for LSQ6 by the forwarding
  \foreach \y in {-0.45,-0.9,-1.35} {
    \draw[thick] (2.83,\y) -- (2.92,\y-0.09) -- (3.08,\y+0.1);
  }
  \draw[thick, ln-accent] (2.83,-2.25) -- (2.92,-2.34) -- (3.08,-2.15);
  % forwarded value in the searching load
  \node[text=ln-accent, font=\footnotesize\bfseries\sffamily] at (2.25,-2.25) {17};
  % pointers
  \draw[->, thick, ln-accent] (-1.45,-0.45) -- (-0.85,-0.45);
  \node[font=\scriptsize\sffamily, text=ln-accent, anchor=east] at (-1.45,-0.45)
    {\iflangja{コミット}{commit}};
  \draw[->, thick, ln-accent] (-1.45,-3.15) -- (-0.85,-3.15);
  \node[font=\scriptsize\sffamily, text=ln-accent, anchor=east] at (-1.45,-3.15)
    {\iflangja{発行}{issue}};
  % search of the load in LSQ6
  \draw[thick, ln-accent] (3.25,-2.25) -- (3.6,-2.25) -- (3.6,-0.45);
  \draw[->, thick, ln-accent] (3.6,-0.45) -- (3.25,-0.45);
  \foreach \y in {-0.9,-1.35,-1.8} { \draw[ln-accent] (3.52,\y) -- (3.68,\y); }
  \node[ann, text=ln-accent] at (3.8,-2.25)
    {\iflangja{アドレス 0x100 のロード：前方を探索}{load at 0x100: searches earlier entries}};
  \node[ann, text=ln-muted] at (3.8,-1.8) {\iflangja{ロード：無視}{load: skipped}};
  \node[ann, text=ln-muted] at (3.8,-1.35) {\iflangja{ストア，0x108：不一致}{store, 0x108: no match}};
  \node[ann, text=ln-muted] at (3.8,-0.9) {\iflangja{ロード：無視}{load: skipped}};
  \node[ann, text=ln-accent] at (3.8,-0.45)
    {\iflangja{ストア，0x100：一致，値 17 を転送}{store, 0x100: match, value 17 forwarded}};
  \node[ann, text=ln-muted] at (3.8,-2.7)
    {\iflangja{後続のストア：探索しない}{later store: not examined}};
  % program order
  \draw[->, ln-muted] (-2.55,0.2) -- (-2.55,-3.35);
  \node[lab, rotate=90] at (-2.8,-1.57) {\iflangja{プログラム順}{program order}};
\end{tikzpicture}

  \caption{8 エントリのロードストアキュー．LSQ2 から LSQ7 までが使用中で
    ある．LSQ6 のロードはアドレス \texttt{0x100} を計算した．その探索は
    ロードと \texttt{0x108} へのストアを飛ばし，\texttt{0x100} への最も
    近い先行ストアから値 17 を取る．}
  \label{fig:l09-lsq}
\end{figure}

ロードはアドレスを計算すると，それを自身の LSQ エントリに記録し，自分の
エントリからコミットポインタに向かって先行するエントリを探索する．位置を
読むことはその内容を変えないので，ロードは飛ばされ，ストアだけが問題に
なる．探索は，同じアドレスを持つ最初のストアで，あるいはコミットポインタ
で終わる．
\begin{enumerate}
  \item 同じアドレスのストアが見つかれば，この最も近い一致するストアが，
    プログラム順でロードの前にその位置へ書き込む最後の命令なので，その値
    こそロードが読むべき値である．ロードはそのストアの LSQ エントリから
    値を複写し，ストアがまだ値を記録していなければそれを待ち，メモリには
    まったくアクセスしない．このように未コミットのストアの値をロードに
    与えることを，\term[すとあろーどふぉわーでぃんぐ]{ストアロードフォワーディング}%
    [store-to-load forwarding]という．
  \item 一致するストアもアドレスの分からないストアも見つからないまま
    コミットポインタに達すれば，ロードの前にその位置へ書き込む未コミット
    のストアは存在せず，メモリがすでに正しい値を保持しているので，ロード
    はメモリを読む．\caution{第3章のフォワーディングとは別物である．
    あちらはパイプラインステージの間で結果をバイパスするが，ここでは値を
    一方の命令の LSQ エントリから他方へ複写する．}
\end{enumerate}
探索が終わる前に，まだアドレスの計算されていないストアを通過した場合，
そのストアはロードの読む位置に書き込むかもしれないし書き込まないかも
しれず，プロセッサにはまだどちらとも分からない．これが難しい場合であり，
次節ではその扱い方を示す．

\section[ロードストアのアウトオブオーダ実行とインオーダ実行]%
  {ロードストアのアウトオブオーダ実行と\\インオーダ実行}
\label{sec:l09-policies}

アドレスの分からないストアを探索が通過してしまうロードに対しては，3つの
選択肢がある．\source{Lesson 9，動画 5--8；H\&P §3.6．}
\begin{enumerate}
  \item 第7章と同様に，ロードとストアを互いにプログラム順に実行する．
  \item 先行するすべてのストアがアドレスを計算するまでロードを待たせ，
    その探索がアドレスの分からないストアに出会わないようにする．
  \item ロードをそのまま実行し，誤った値は後で検出して訂正する．
    \margin{分岐予測（第4章）と同じ型である．予測して先に進み，後で検査
    し，誤っていればそれ以降を取り消す．}
\end{enumerate}
最初の2つはロードを待たせ，3つ目は待たせない代わりに検査を必要とする．
ストアのアドレスはベースレジスタが得られるまで分からず，そのレジスタを
遅い命令が計算する場合には長くかかりうる．次のプログラムでは，ストアの
アドレスは乗算に依存するのに対し，ロードのアドレスはほとんど直ちに
得られる．
\begin{lstlisting}[language={[x86masm]Assembler}, numbers=none]
I1:  MUL R1, R2, R3     ; R1 <- R2 * R3
I2:  ADD R1, R1, R4     ; R1 <- R1 + R4（アドレス）
I3:  SW  R5, 0(R1)      ; MEM[R1] <- R5
I4:  ADD R6, R6, R7     ; R6 <- R6 + R7（アドレス）
I5:  LW  R8, 0(R6)      ; R8 <- MEM[R6]
I6:  SUB R9, R8, R5     ; R9 <- R8 - R5
\end{lstlisting}
タイミングの規約は \cref{sec:l08-example}のものを用いる．MUL は 4
サイクル，ADD，SUB，ロード，ストアは 1 サイクルを要し，その 1 サイクルで
ロードはアドレスを計算して値を得，ストアはアドレスを計算する．値 R5 は
最初から利用できる．I1 は結果をサイクル~6 に，I2 はサイクル~8 に書き込む
ので，ストア I3 が実行され，そのアドレスが分かるのはサイクル~9 である．I4 は
結果をサイクル~6 に書き込むので，ロード I5 はサイクル~7 に実行できる
（\cref{fig:l09-policies}）．

\begin{figure}[htbp]
  \centering
  % Out-of-order versus in-order load-store execution for the program of
% section "Out-of-order and in-order load-store execution".  The store I3
% computes its address only in cycle 9; the load I5 could execute in cycle 7.
\begin{tikzpicture}[x=0.52cm,y=0.5cm, font=\scriptsize\sffamily, >={Stealth[length=1.6mm]},
  c/.style={minimum width=0.52cm, minimum height=0.5cm, inner sep=0pt, draw=black},
  iss/.style={c, fill=white},
  ex/.style={c, fill=ln-accent!25},
  sto/.style={c, fill=ln-caution!20},
  wr/.style={c, fill=ln-accent, text=white, font=\scriptsize\bfseries\sffamily},
  hd/.style={font=\scriptsize\bfseries\sffamily, anchor=west}]
  % cycle axes and grids of both panels
  \foreach \ya/\yb/\yc in {0/-0.4/-6.5, -8/-8.4/-11.5} {
    \foreach \c in {1,...,13} { \node[text=ln-muted] at (\c,\ya) {\c}; }
    \node[anchor=east, text=ln-muted] at (0.4,\ya) {\iflangja{サイクル}{cycle}};
    \foreach \c in {1,...,14} { \draw[ln-rule, very thin] (\c-0.5,\yb) -- (\c-0.5,\yc); }
  }
  % ------------------------------------------------ (a) out-of-order
  \node[hd] at (-6.2,0.9) {\iflangja{(a) アウトオブオーダのロード／ストア実行}%
                                     {(a) out-of-order load-store execution}};
  \foreach \y/\lab/\i/\ea/\eb/\w/\st in {
      -1/{I1\enskip\texttt{MUL R1,R2,R3}}/1/2/5/6/ex,
      -2/{I2\enskip\texttt{ADD R1,R1,R4}}/2/7/7/8/ex,
      -3/{I3\enskip\texttt{SW\ \ R5,0(R1)}}/3/9/9/0/sto,
      -4/{I4\enskip\texttt{ADD R6,R6,R7}}/4/5/5/6/ex,
      -5/{I5\enskip\texttt{LW\ \ R8,0(R6)}}/5/7/7/8/ex,
      -6/{I6\enskip\texttt{SUB R9,R8,R5}}/6/9/9/10/ex,
      -9/{I3\enskip\texttt{SW\ \ R5,0(R1)}}/3/9/9/0/sto,
      -10/{I5\enskip\texttt{LW\ \ R8,0(R6)}}/5/10/10/11/ex,
      -11/{I6\enskip\texttt{SUB R9,R8,R5}}/6/12/12/13/ex} {
    \node[anchor=east] at (0.4,\y) {\lab};
    \node[iss] at (\i,\y) {I};
    \pgfmathtruncatemacro{\lnwa}{\i+1}
    \pgfmathtruncatemacro{\lnwb}{\ea-1}
    \ifnum\lnwa<\ea
      \draw[ln-muted, thick] (\lnwa-0.5,\y) -- (\lnwb+0.5,\y);
    \fi
    \foreach \e in {\ea,...,\eb} { \node[\st] at (\e,\y) {E}; }
    \ifnum\w>0 \node[wr] at (\w,\y) {W}; \fi
  }
  % check of the store
  \draw[->, ln-caution, thick] (9.28,-3.25) -- (9.28,-4.3) -- (7.3,-4.3) -- (7.3,-4.72);
  \node[anchor=west, align=left, text=ln-caution, fill=white, inner sep=1pt] at (10.6,-4.5)
    {\iflangja{I3 がアドレスを計算し，\\実行済みの I5 を検査}%
              {I3 computes its address\\and checks I5}};
  % ------------------------------------------------ (b) in-order
  \node[hd] at (-6.2,-7.1) {\iflangja{(b) インオーダのロード／ストア実行}%
                                      {(b) in-order load-store execution}};
  \draw[->, ln-accent, thick] (9.3,-9.25) -- (9.7,-9.75);
  % ------------------------------------------------ legend
  \begin{scope}[shift={(0,-12.6)}]
    \node[iss] at (-5.6,0) {I}; \node[anchor=west] at (-5.1,0) {\iflangja{発行}{issue}};
    \draw[ln-muted, thick] (-2.7,0) -- (-1.7,0);
    \node[anchor=west] at (-1.6,0) {\iflangja{待機}{waiting}};
    \node[ex] at (1.6,0) {E}; \node[anchor=west] at (2.1,0) {\iflangja{実行}{execute}};
    \node[sto] at (5.0,0) {E}; \node[anchor=west] at (5.5,0) {\iflangja{ストア実行}{store executes}};
    \node[wr] at (10.0,0) {W}; \node[anchor=west] at (10.5,0) {\iflangja{結果書き込み}{write result}};
  \end{scope}
\end{tikzpicture}

  \caption{\cref{sec:l09-policies}のプログラムを2つの方針で実行した場合．
    (a) I5 はストア I3 がアドレスを知る前に実行され，I3 はサイクル~9 に
    I5 を検査する．(b) I5 は I3 がアドレスを計算するまで待つ．I1，I2，I4
    は (a) と同じである．}
  \label{fig:l09-policies}
\end{figure}

\subsection{ロードストアのアウトオブオーダ実行}

\begin{definition}[ロードストアのアウトオブオーダ実行]
  \label{def:l09-ooo}
  \term[ろーどすとあのあうとおぶおーだじっこう]{ロードストアのアウトオブオーダ実行}%
  [out-of-order load-store execution]（選択肢~3）
  では，ロードは，先行するストアがアドレスを計算しているかどうかに
  かかわらず，自身のアドレスが計算でき次第実行される．その LSQ 探索は，
  アドレスの分からないストアを，アドレスが異なるものとみなして通過する．
\end{definition}

例では I5 はサイクル~7 に，その値を使う I6 はサイクル~9 に実行される．
その代わり，すべてのストアは自身のアドレスを，自分を追い越したロードと
照合しなければならない．\tip{ロードは自分より前のストアを探し，ストアは
自分より後のロードを探す．どちらの探索も自分と同じ種類は飛ばす．}I3 は
サイクル~9 にアドレスを計算すると，後続の LSQ エントリの中から，同じ
アドレスを持ち C ビットの立ったロード，すなわち
すでに値を得たロードを探す．
\begin{itemize}
  \item I5 のアドレスが異なれば何も起こらない．I5 は正しい値を得ており，
    I6 はサイクル~10 に結果を書き込む．
  \item アドレスが等しければ，I5 は，プログラム順では I3 が上書きする値を
    得てしまっているので，その値を直接または間接に使ったすべての命令が
    誤った結果を計算したことになる．プロセッサは I5 とそれ以降のすべての
    命令を取り消し，I5 のフェッチからやり直して再実行する．I3 がそれらに
    先行するので，いずれもコミットしておらず，ROB は予測を外した経路の
    命令と同じようにそれらを破棄する（\cref{sec:l08-recovery}）．
    \margin{取り消しにはパイプラインの再起動の費用がかかる．Skylake 級の
    コアで 15--20 サイクルであり，分岐予測ミスと同程度である（Agner Fog，
    マイクロアーキテクチャマニュアル）．}I5 が
    再び実行されると，その探索は I3 を見つけ，フォワーディングによって
    値を受け取る．
\end{itemize}

\subsection{ロードストアのインオーダ実行}

\begin{definition}[ロードストアのインオーダ実行]
  \label{def:l09-in-order}
  \term[ろーどすとあのいんおーだじっこう]{ロードストアのインオーダ実行}%
  [in-order load-store execution]（選択肢~1）では，ロードとストアは互いに
  プログラム順に実行される．すなわち，ロードやストアは，先行するすべての
  ロードとストアが実行されるまで待つ．ここでストアは，アドレスが分かった
  時点で実行されたとみなす．他の命令は依然としてアウトオブオーダに実行
  される．
\end{definition}

例では I5 は I3 を待ってサイクル~10 に実行され，I6 はサイクル~12 に実行
されてサイクル~13 に結果を書き込む．どのロードも誤った値を得ることはなく，
再実行も起こらない．しかし I5 は，そのアドレスが I3 のアドレスと異なる
場合にも 3 サイクル待ち，そのときこの待ちは純粋な損失である．ここでは
選択肢~2 も同じタイミングになる．I5 に先行するロードがないからである．
一般には選択肢~2 のほうが待ちは短い．ロードは先行するロードには依存
しないからである．

\subsection{比較}

ロードストアのアウトオブオーダ実行は，ロードが，アドレスのまだ分からない
先行ストアの書き込む位置を読まない限りサイクルを稼ぎ，読む場合にはサイクル
を失う．ロードとそれ以降の命令を取り消して再実行するコストは，待った場合
よりも大きいからである．選択肢 1 と 2 は，アドレスの分からないストアに
続くすべてのロードを，両者が衝突するかどうかにかかわらず遅らせる．\sidenote{実際のコアは，
一律に決めるのではなくロードごとに決める．そのロードが衝突しそうかどうかを
予測器が判定するのである．Alpha 21264 は，過去に取り消しを起こした PC を
記録した表に載るロードだけを待たせ，その表を定期的に消去して，ロードが再び
投機的になれるようにする．ストアセット \cite{chrysos1998} は，各ロードを，
それと衝突したことのあるストアに対応付ける．Intel はこの仕組みを memory
disambiguation と呼び，2006 年から搭載している．}%
実際のプログラムでは，そのようなストアの書き込む内容を本当に読むロードはごく
わずかなので，ロードストアのアウトオブオーダ実行のほうが性能がよく，
高性能プロセッサが採用しているのはこの方針である．

\begin{keypoint}
  ロードストアのインオーダ実行はロードを再実行することがない代わりに，
  通常は依存しないストアをロードに待たせる．ロードストアのアウトオブ
  オーダ実行はロードを直ちに実行させ，その正しさは，各ストアがアドレスを
  計算するときに完了済みの後続ロードを検査することに依存する．
\end{keypoint}

\section{ストアロードフォワーディング}
\label{sec:l09-forwarding}

ストアロードフォワーディングは，各ロードと，その値を読むべきストアとを
対応付ける．ロードストアのアウトオブオーダ実行を行うプロセッサは，この
対応を両側から用いる．\source{Lesson 9，動画 9；H\&P 第3章．}
\begin{description}
  \item[ロードにとって] 値は，同じアドレスを持つ最も近い先行ストア，
    すなわちプログラム順でロードの前にその位置へ書き込む最後のストアから
    来る．同じアドレスへのより古いストアは該当しない．その値はより近い
    ストアに上書きされるからである．LSQ にこのアドレスの先行ストアが
    なければ，値はメモリから来る．\caution{探索はロードから遡って行う．
    コミットポインタから順に前向きに探索すると，一致するストアのうち最も
    古いものが見つかってしまう．}
  \item[ストアにとって] 値は，同じアドレスを持つ後続のすべてのロードに，
    そのアドレスへの次のストアの手前まで届く．そのストアより後のロードは，
    代わりにそのストアの値を受け取る．\margin{ハードウェアがフォワードでき
    るのは，1つのストアの中に完全に収まるロードだけである．一部だけ重なる
    場合，ロードはストアが L1 に届くまで待つことになり，12 サイクルほど
    かかる（Agner Fog）．}
\end{description}
ロードは探索（\cref{sec:l09-lsq}）で第1の規則を，ストアは完了済みの後続
ロードの検査（\cref{sec:l09-policies}）で第2の規則を適用する．
\cref{fig:l09-forwarding} は，2つのアドレスに対する7つのロードとストアに
ついてこの対応を示す．

\begin{figure}[htbp]
  \centering
  % Store-to-load forwarding for a sequence of seven loads and stores: each load
% receives the value of the nearest earlier store with the same address, or
% the value in memory if there is none.  A store's value reaches the later
% loads with its address up to the next store with that address.
\begin{tikzpicture}[x=1cm,y=1cm, font=\footnotesize\sffamily, >={Stealth[length=1.8mm]},
  hd/.style={font=\footnotesize\bfseries\sffamily},
  row/.style={draw, minimum height=0.45cm, inner sep=0pt},
  lab/.style={font=\scriptsize\sffamily, text=ln-muted}]
  % column headings
  \node[hd] at (0.35,0.45) {\iflangja{命令}{Entry}};
  \node[hd] at (1.05,0.45) {L/S};
  \node[hd] at (1.95,0.45) {Addr};
  \node[hd] at (2.9,0.45) {Value};
  \foreach \r/\n/\t/\a/\v/\f in {1/S1/S/0x40/5/ln-accent!12, 2/L2/L/0x48/30/ln-tip!12,
      3/S3/S/0x48/6/ln-tip!12, 4/L4/L/0x40/5/ln-accent!12, 5/S5/S/0x40/9/ln-accent!12,
      6/L6/L/0x48/6/ln-tip!12, 7/L7/L/0x40/9/ln-accent!12} {
    \pgfmathsetmacro{\y}{-(\r-1)*0.45}
    \node[row, fill=ln-box, minimum width=0.7cm, text=ln-muted] at (0.35,\y) {\n};
    \node[row, fill=\f, minimum width=0.7cm] at (1.05,\y) {\t};
    \node[row, fill=\f, minimum width=1.1cm, font=\scriptsize\ttfamily] at (1.95,\y) {\a};
    \node[row, fill=\f, minimum width=0.8cm] at (2.9,\y) {\v};
  }
  % address 0x40 (blue): S1 -> L4, S5 -> L7 (right side)
  \draw[thick, ln-accent] (3.3,0.0) -- (3.75,0.0) -- (3.75,-1.35);
  \draw[->, thick, ln-accent] (3.75,-1.35) -- (3.3,-1.35);
  \draw[thick, ln-accent] (3.3,-1.8) -- (3.75,-1.8) -- (3.75,-2.7);
  \draw[->, thick, ln-accent] (3.75,-2.7) -- (3.3,-2.7);
  \node[font=\scriptsize\sffamily, text=ln-accent, anchor=west, align=left] at (3.95,-0.675)
    {\iflangja{S1 の値：\\S5 の手前まで}{value of S1:\\up to S5}};
  \node[font=\scriptsize\sffamily, text=ln-accent, anchor=west, align=left] at (3.95,-2.25)
    {\iflangja{S5 の値}{value of S5}};
  % address 0x48 (green): memory -> L2, S3 -> L6 (left side)
  \node[draw, fill=white, font=\scriptsize\sffamily, inner sep=2pt] (mem) at (-1.35,-0.45)
    {\iflangja{メモリ}{memory}};
  \draw[->, thick, ln-tip!80!black] (mem.east) -- (0.0,-0.45);
  \draw[thick, ln-tip!80!black] (0.0,-0.9) -- (-0.45,-0.9) -- (-0.45,-2.25);
  \draw[->, thick, ln-tip!80!black] (-0.45,-2.25) -- (0.0,-2.25);
  \node[font=\scriptsize\sffamily, text=ln-tip!80!black, anchor=east, align=right] at (-0.6,-1.575)
    {\iflangja{S3 の値}{value of S3}};
  % program order
  \draw[->, ln-muted] (5.75,0.2) -- (5.75,-2.9);
  \node[lab, rotate=-90] at (6.0,-1.35) {\iflangja{プログラム順}{program order}};
\end{tikzpicture}

  \caption{\texttt{0x40} と \texttt{0x48} への7つのロードとストアに対する
    ストアロードフォワーディング．値は逐次実行の場合のものである．L4 は
    S1 の値を受け取り，その後ろにある S5 の値は受け取らない．L7 は S1 の
    値を上書きする S5 の値を受け取る．L2 は S3 に先行するのでメモリを
    読む．}
  \label{fig:l09-forwarding}
\end{figure}

\section{ロードストアキューの例}
\label{sec:l09-example}

ここで，\cref{fig:l09-forwarding} のロードとストア（\cref{ex:l09-lsq}）を，
ロードストアのアウトオブオーダ実行のもとでトレースする．
\source{Lesson 9，動画 10；H\&P §3.6．}

\begin{example}
  \label{ex:l09-lsq}
  あるプログラムのロードとストアは，プログラム順に次のとおりである．
  \begin{lstlisting}[language={[x86masm]Assembler}, numbers=none]
S1:  SW  R1, 0(R2)     ; MEM[0x40] <- R1
L2:  LW  R3, 8(R2)     ; R3 <- MEM[0x48]
S3:  SW  R4, 8(R10)    ; MEM[R10+8] <- R4
L4:  LW  R5, 0(R2)     ; R5 <- MEM[0x40]
S5:  SW  R6, 0(R7)     ; MEM[R7] <- R6
L6:  LW  R8, 8(R2)     ; R8 <- MEM[0x48]
L7:  LW  R9, 0(R2)     ; R9 <- MEM[0x40]
\end{lstlisting}
  ここで $\text{R2} = \texttt{0x40}$，$\text{R1}=5$，$\text{R4}=6$，
  $\text{R6}=9$ は最初から利用できる．$\text{R10} = \texttt{0x40}$ と
  $\text{R7}=\texttt{0x40}$ は S3 と S5 に先行する遅い命令が計算するので，
  この2つのストアがアドレスを知るのは遅く，その順序は
  \cref{tab:l09-trace} のとおりである．メモリは最初 \texttt{0x40} に 20
  を，\texttt{0x48} に 30 を保持している．\tip{トレースを検算するには，
  まず逐次実行での値を書き出せばよい．それと異なる値で終わるロードが
  あれば，再実行が抜けている．}
\end{example}

\begin{table}[htbp]
  \centering\small
  \caption{\cref{ex:l09-lsq} の実行．各ステップで1つのエントリがアドレス
    を計算する．ロードは先行するストアを探索し，ストアは C ビットの立った
    後続ロードを検査する．}
  \label{tab:l09-trace}
  \begin{tabular}{@{}clcll@{}}
    \toprule
    ステップ & エントリ & Addr & 探索または検査 & 結果 \\
    \midrule
    1 & S1 & \texttt{0x40} & 完了した後続ロードなし & 値 5 を記録 \\
    2 & L2 & \texttt{0x48} & S1 は相違 & メモリから 30 \\
    3 & L4 & \texttt{0x40} & S3 は不明，S1 が一致 & S1 から 5 \\
    4 & S3 & \texttt{0x48} & L4 は相違 & 衝突なし \\
    5 & L6 & \texttt{0x48} & S5 は不明，S3 が一致 & S3 から 6 \\
    6 & L7 & \texttt{0x40} & S5 は不明，S3 は相違，S1 が一致 & S1 から 5 \\
    7 & S5 & \texttt{0x40} & L6 は相違，L7 が一致 & L7 を再実行 \\
    8 & L7 & \texttt{0x40} & S5 が一致 & S5 から 9 \\
    \bottomrule
  \end{tabular}
\end{table}

L4 と L6 は，アドレスのまだ分からないストアを通過するが，その後の S3 と
S5 の検査では，これらとの衝突は見つからない．L7 も S5 を通過するが，
ステップ~7 で S5 が \texttt{0x40} を計算すると，このアドレスで完了して
いる L7 が見つかり，しかも両者の間に \texttt{0x40} へのストアはない．
L7 が S1 から得た値 5 は，プログラム順では S5 に上書きされるので，L7 と
それ以降のすべての命令が取り消されて再実行され，ステップ~8 で L7 は S5
から 9 を受け取る．

エントリがプログラム順にコミットした後，メモリは \texttt{0x40} に 9 を，
\texttt{0x48} に 6 を保持し，ロードは $\text{R3}=30$，$\text{R5}=5$，
$\text{R8}=6$，$\text{R9}=9$ を返している．これは逐次実行の結果と同じで
ある．ロードストアのインオーダ実行であれば，L4 は S3 を，L6 と L7 は S5
を待ち，二度実行される命令はなかったはずである．

\section{LSQ と ROB とリザベーションステーション}
\label{sec:l09-rob-rs}

ロードとストアも他の命令と同じステップを経るが，LSQ が各ステップに作業を
加える（\cref{tab:l09-steps}）．\source{Lesson 9，動画 11--14；H\&P
§3.6．}ロードやストアが発行されるのは，ROB エントリと LSQ エントリの両方
が空いているときだけである．LSQ エントリは，他の命令が ROB エントリと
ともに必要とするリザベーションステーションの代わりとなる．\margin{実際の
規模：Skylake は 224 エントリの ROB に対して 72 個のロードと 56 個のストアを，
Ice Lake は 352 に対して 128 個と 72 個を実行中に保持できる（Intel 最適化
マニュアル）．}ロードは，算術
演算の結果とまったく同じように，自身の ROB エントリの名前のもとで値を CDB
にブロードキャストする．ストアはレジスタの値を生成しないので何も
ブロードキャストせず，その ROB エントリは，ストアがアドレスと値の両方を
記録した時点で完了となる．

\begin{table}[htbp]
  \centering\small
  \caption{LSQ を持つプロセッサにおけるロードとストアのステップ．}
  \label{tab:l09-steps}
  \begin{tabular}{@{}>{\raggedright\arraybackslash}p{0.15\linewidth}>{\raggedright\arraybackslash}p{0.38\linewidth}>{\raggedright\arraybackslash}p{0.38\linewidth}@{}}
    \toprule
    ステップ & ロード & ストア \\
    \midrule
    発行 & 空いている ROB エントリと LSQ エントリを必要とする（RS は不要）．
      ベースレジスタを RAT で引き，宛先レジスタの RAT エントリをそのロード
      の ROB エントリに設定する & 空いている ROB エントリと LSQ エントリを
      必要とする（RS は不要）．ベースレジスタと値レジスタを RAT で引く．
      RAT エントリは変化しない \\
    \addlinespace
    実行 & アドレスを計算して LSQ に記録し，先行するストアを探索し，
      フォワードされた値を複写するかメモリを読み，C を立てる
      & ベースレジスタが得られたときアドレスを記録し，C の立った後続ロード
      を検査する．値は得られたときに記録する．両方を記録した時点で C を
      立てる \\
    \addlinespace
    結果書き込み & 値を CDB にブロードキャストし，ROB エントリがそれを
      受け取る & ブロードキャストするものはない \\
    \addlinespace
    コミット & 値をレジスタファイルに書き込み，RAT エントリがまだこの ROB
      エントリを指していればそれを消す．ROB エントリと LSQ エントリを
      解放する & 値をメモリに書き込む．ROB エントリと LSQ エントリを解放
      する \\
    \bottomrule
  \end{tabular}
\end{table}

% 最後の段落を章のまとめと同じページに置く（まとめだけのページを避ける）
\needspace{21\baselineskip}
どちらのキューもプログラム順なので，ROB のコミットポインタにあるロード
またはストアは LSQ の最も古いエントリでもあり，2つのエントリは同時に解放
される．\index{こみっと@コミット}ストアが LSQ を離れるのは，その値が
メモリに書き込まれるのとちょうど同じ時点なので，後続のロードは，ロードが
探す2つの場所のいずれかで必ずその値を見つける．コミット前はストアの LSQ
エントリに，コミット後はメモリにある．\margin{マルチコアではもう1つ規則が
要る．すでに読み終えたロードも，コミットより前に他のコアがそのアドレスへ
書き込んだ場合には取り消さなければならない（第21章）．}

\begin{keypoint}[章のまとめ]
  ロードとストアは，実行されてはじめて分かるメモリアドレスを介して互いに
  依存し，メモリはリネーミングされない．そのためストアは，コミット時に
  プログラム順でメモリに書き込む．ロードストアキューは，未コミットの
  ロードとストアを，そのアドレス，値，完了ビットとともにプログラム順に
  保持する．ロードは，ストアロードフォワーディングによって，同じアドレス
  を持つ最も近い先行ストアから値を取り，そのようなストアがなければメモリ
  から取る．アドレスの分からない先行ストアの扱いには2つの方針がある．
  ロードストアのインオーダ実行ではロードがそれらを待ち，ロードストアの
  アウトオブオーダ実行ではロードが先に進み，後に一致するアドレスを計算
  したストアが，そのロードとそれ以降のすべての命令を再実行させる．
  そのような衝突は稀なので，後者のほうが速い．ロードとストアは発行時に
  ROB エントリと LSQ エントリを必要とし，他の命令は ROB エントリと RS を
  必要とする．
\end{keypoint}

\end{document}
